
\uline{est un
isomorphisme}.

Pour voir qu'il est injectif, on utilise la suite d'homomorphismes canoniques~:
@@M0@@
la dernière flèche venant de ce que @@M1@@ est un sous-foncteur de @@M2@@. Pour vérifier que l'homomorphisme composé est l'identité, on remarque qu'un morphisme de foncteurs~: @@M3@@ est déterminé par sa valeur pour les objets de @@M4@@ dont la filtration est grossière~; en effet, soit @@M5@@, et @@M6@@ l'algèbre @@M7@@ munie de la filtration grossière~; le morphisme canonique @@M8@@ définit un diagramme commutatif~:
@@M9@@
montrant que @@M10@@ est la restriction de @@M11@@ à @@M12@@.

On considère ensuite l'homomorphisme composé~:
@@M13@@
et il est clair que c'est l'identité, d'où le résultat annoncé.

\textbf{Théorème 1.10.} \uline{Soit} @@M14@@ \uline{un foncteur sur la catégorie} @@M15@@ \uline{des} @@M16@@-\uline{algèbres}, \uline{à valeurs dans la catégorie des groupes abéliens filtrés séparés complets}, \uline{et commutant aux invariants sous un groupe fini}. \uline{Alors l'application de} @@M17@@ \uline{dans} @@M18@@ \uline{qui}, \uline{à un morphisme de foncteurs}, \uline{associe}
